\section{Method}
\label{sec:method}

\begin{figure*}[t]
\centering
\resizebox{\textwidth}{!}{%
\begin{tikzpicture}[
  font=\small,
  box/.style={draw, rounded corners=2pt, minimum height=0.9cm, align=center, inner sep=4pt},
  fz/.style={box, fill=frozen!15, draw=frozen, thick},
  tr/.style={box, fill=trained!15, draw=trained, thick},
  tok/.style={draw, minimum width=0.38cm, minimum height=0.42cm, inner sep=0pt},
  arr/.style={-{Stealth[length=2.2mm]}, thick}]
  % inputs and the two streams
  \node[box, fill=gray!10, minimum width=1.9cm] (x) at (0,0) {Meeting\\audio $\bm{x}$};
  \node[fz, minimum width=2.7cm] (enc) at (3.2,1.0) {Audio encoder\\(\S\ref{sec:method-task})};
  \node[fz, minimum width=2.7cm] (dia) at (3.2,-1.0) {Streaming diarizer\\$S{=}8$ slots};
  \node[tr, minimum width=2.3cm] (proj) at (6.5,-1.0) {Speaker projector\\(App.~\ref{app:model})};
  \draw[arr] (x.east) -- ++(0.3,0) |- (enc.west);
  \draw[arr] (x.east) -- ++(0.3,0) |- (dia.west);
  \draw[arr] (dia) -- node[above]{$\bm{P}$} (proj);
  % interleaved sequence
  \foreach \i/\c in {0/frozen,1/frozen,2/frozen,3/frozen,4/frozen,5/spkc,6/frozen,7/frozen,8/frozen,9/frozen,10/frozen,11/spkc} {
    \node[tok, fill=\c!35] (t\i) at (8.9+\i*0.4,0) {};
  }
  \node (dots) at (8.9+12*0.4+0.05,0) {$\cdots$};
  \draw[decorate, decoration={brace, amplitude=4pt}] (t0.north west) -- node[above=4pt, font=\small]{$\bm{A}_1$: $K{=}5$ acoustic tokens} (t4.north east);
  \node[below=2pt of t5, font=\small] {$\bm{s}_1$};
  \node[below=2pt of t11, font=\small] {$\bm{s}_2$};
  \node[font=\small] at (11.2,-1.25) {interleaved input $\bm{E}$ (Eq.~\ref{eq:interleave})};
  \draw[arr] (enc.east) -| node[pos=0.3, above, font=\small]{acoustic tokens $\bm{a}_i$} (t8.north);
  \draw[arr] (proj.east) -- node[above, font=\small]{$\bm{s}_j$} ++(0.75,0) |- (t0.west);
  % LLM
  \node[tr, minimum height=1.9cm, minimum width=2.0cm] (llm) at (15.6,0) {Thinker LLM\\+ LoRA\\(Eq.~\ref{eq:loss})};
  \draw[arr] (dots.east) -- (llm.west);
  % output segments (second row)
  \node[box, fill=white, anchor=west] (o1) at (6.2,-2.6)
    {$\spk{2}\,\tim{5}$\ \zh{我想知道会议的时间}\ $\tim{48}$};
  \node[box, fill=white, anchor=west] (o2) at (6.2,-3.7)
    {$\spk{3}\,\tim{42}$\ \zh{对，没错。}\ $\tim{56}$};
  \draw[arr] (llm.south) -- ++(0,-0.85) -| ([xshift=-0.3cm]o1.west) -- (o1.west);
  \draw[arr] ([xshift=-0.3cm]o1.west) |- (o2.west);
  \node[font=\small, anchor=west] at (6.2,-4.45) {who, when, what (Eq.~\ref{eq:serial})};
  % talker
  \node[tr, minimum width=2.4cm] (talk) at (14.9,-3.15) {Talker, voice $\bm{v}_k$\\playback alignment\\(\S\ref{sec:method-plem})};
  \draw[arr] (o1.east) -- ++(0.25,0) |- (talk.west);
  \draw[arr] (o2.east -| o1.east) ++(0,0) -- ++(0.25,0) |- (talk.west);
  \draw[thick] (o2.east) -- (o2.east -| o1.east);
  \node[box, fill=gray!10, minimum width=1.7cm] (y) at (17.6,-3.15) {Translated\\speech};
  \draw[arr] (talk) -- (y);
  % voice prompts
  \draw[arr, dashed, spkc!80!black] (dia.south) -- ++(0,-3.75) -| node[pos=0.25, below, font=\small]{voice prompt $\bm{v}_k$: solo regions of slot $k$ (Eq.~\ref{eq:alone})} (talk.south);
  % streaming loop
  \draw[arr, dotted, trained!80!black] (llm.north) -- ++(0,0.5) -- node[above, font=\small]{streaming: read every $\Delta$\,s, write or $\wait$ on deadlines (Eq.~\ref{eq:write})} ++(-7.4,0);
\end{tikzpicture}}
\caption{\sys{}. A frozen audio encoder turns the meeting into acoustic tokens (what is said); a frozen streaming
diarizer gives slot activities $\bm{P}$ that a trained projector turns into speaker tokens $\bm{s}_j$ (who says it).
One speaker token follows every block of $K{=}5$ acoustic tokens. The Thinker writes speaker-attributed, time-stamped
translation segments, offline or read by read, and a Talker speaks each segment in the voice of its speaker slot.
Blue: frozen; orange: trained. Dotted: the streaming loop, which writes on deadlines; the Talker's playback is trained by playback alignment. The Thinker's gradients reach only the projector, the new token rows and its LoRA
adapters; the Talker is fine-tuned separately on the playback timeline (\S\ref{sec:method-plem}) with the Thinker
fixed. The diarizer and the encoder are never trained.}
\label{fig:arch}
\end{figure*}


\subsection{Task and input: who said what}
\label{sec:method-task}
\label{sec:method-audio}
\label{sec:method-spk}
\label{sec:method-fuse}

\paragraph{Task.} From a single-channel recording $\bm{x}$ of $T$ seconds in which an unknown number of speakers talk,
sometimes at once, the system outputs segments
$\mathcal{Y}=\big((k_n,a_n,b_n,\bm{y}^{(n)})\big)_{n=1}^{N_\mathcal{Y}}$, where $k_n\in\{1,\dots,S\}$ is the speaker label (arbitrary, but one person keeps one label for the whole meeting),
$a_n<b_n$ the start and end of the source speech and $\bm{y}^{(n)}$ its translation; in speech output $\bm{y}^{(n)}$ is
spoken in the voice of $k_n$. In streaming mode segments are produced while $\bm{x}$ is recorded.

\paragraph{Two streams.} The frozen audio encoder of Qwen3-Omni, whose LLM (the ``Thinker'') already translates
single-speaker speech~\citep{xu2025qwen3omni}, maps the recording to acoustic tokens
$\bm{A}=\mathrm{Enc}(\bm{x})=[\bm{a}_1;\dots;\bm{a}_N]\in\R^{N\times d}$ (one per 80\,ms, centred at $\bar{t}_i$; $d=2048$), which carry no memory of who spoke earlier. A frozen Streaming Sortformer~\citep{medennikov2025streaming}, run in its
streaming mode, keeps every voice in the slot it got when it first appeared and gives slot activities
$\bm{P}=\mathrm{Diar}(\bm{x})\in[0,1]^{F\times S}$  for $F$ frames
of 80\,ms (frame $f$ at time $t_f$) and $S=8$ slots; overlapped speech appears as several large entries in one row. We resample $\bm{P}$
to the token centres, $\tilde{\bm{P}}=\mathrm{Interp}(\bm{P};\bar{t}_{1:N})$, and summarise every block of $K=5$ rows by a strided
convolution and a two-layer perceptron into one speaker token $\bm{s}_j\in\R^{d}$ (6.3M parameters; the output layer starts at zero, so training
starts from the pretrained Thinker; Appendix~\ref{app:model}). A streaming frame is emitted once and never revised, so a
reader at time $t$ sees what a run on $\bm{x}_{\le t}$ gives (Appendix~\ref{app:causal}).

\paragraph{Interleaved fusion.} Following G-STAR~\citep{peng2026gstar}, each speaker token follows the audio it describes,
$\bm{E}=[\bm{A}_1;\bm{s}_1;\dots;\bm{A}_M;\bm{s}_M]$ with $\bm{A}_j$ the $j$-th block of $K$ acoustic tokens;
$\bm{E}$ fills the audio placeholder of the chat template after an instruction
$\bm{\iota}$ (Appendix~\ref{app:prompt}).

\subsection{Speaker-attributed, time-stamped output}
\label{sec:method-out}
\label{sec:method-train}

\paragraph{Serialisation.} We add 251 time tokens $\tim{0},\dots,\tim{250}$ (80\,ms steps up to 20\,s) and 8 speaker
tokens $\spk{1},\dots,\spk{8}$ with trainable input and output rows, and write a window's segments speaker first,
ordered by start as in serialized output training~\citep{kanda2020sot},
\begin{equation}
\bm{y}=\mathop{\Vert}_{n=1}^{N_\mathcal{Y}}\big[\spk{k_n}\,\tim{\alpha_n}\,\bm{y}^{(n)}\,\tim{\beta_n}\big],
\label{eq:serial}
\end{equation}
with $\alpha_n,\beta_n$ start and end rounded to 80\,ms. Targets use the diarizer's slots: per meeting, each reference
speaker is mapped to the slot that overlaps most with that speaker's speech over the whole meeting (Hungarian
matching;~\citealp{kuhn1955hungarian}), so the model learns to copy the slot rather than invent a label per window
(Appendix~\ref{app:model}).

\paragraph{Training.} The Thinker $p_\theta(\bm{y}\mid\bm{E},\bm{\iota})=\prod_u p_\theta(y_u\mid\bm{y}_{<u},\bm{E},\bm{\iota})$ is trained with the weighted cross-entropy
\begin{equation}
\mathcal{L}=-\frac{1}{\Omega}\sum_{u}\omega(y_u)\log p_\theta(y_u\mid\bm{y}_{<u},\bm{E},\bm{\iota}),
\label{eq:loss}
\end{equation}
$\Omega=\sum_u\omega(y_u)$, $\omega=2$ for speaker and $1.5$ for time tokens, since one wrong speaker or time token can misattribute a sentence.
Trained are the projector, the new rows and LoRA adapters~\citep{hu2022lora} on all 48 layers (Table~\ref{tab:hyper});
\sys{} (offline) decodes a window at once.

\subsection{Streaming interpretation}
\label{sec:method-simul}

\paragraph{Reading and writing.} The model reads every $\Delta=1$\,s, at $t_r=\min(r\Delta,T)$, $r=1,\dots,R$. It sees
$\bm{x}_{\le t_r}$ and a causal speaker stream: the low-latency diarizer finalises a frame at most $\gamma=1.04$\,s later,
so later frames are zeroed,
\begin{equation}
P^{(r)}_{f,k}=P_{f,k}\;\mathds{1}\big[t_f\le t_r-\gamma\big],
\label{eq:causal}
\end{equation}
and $\bm{E}^{(r)}$ is built from $\bm{x}_{\le t_r}$ and $\bm{P}^{(r)}$. With $\hat{\bm{y}}^{(<r)}$ everything committed
before read $r$, the model continues greedily from $(\hat{\bm{y}}^{(<r)},\bm{E}^{(r)},\bm{\iota})$ until a stop
token, $\eta_r=\wait$ for $r<R$ and $\eta_R=\eos$; what precedes it is committed and never revised.
Eq.~\eqref{eq:causal} uses the bound $\gamma$; the actual release time of each frame (Appendix~\ref{app:chunk}) is used
to time commits for scoring.

\paragraph{Deadline targets.} Each reference translation is split at punctuation into chunks, and each chunk is timed by
the source it covers (word alignments, or the same fraction of the utterance where none exist; Appendix~\ref{app:model}).
A chunk $j$ ending at $e_j$ becomes due one interpreting lag $\epsilon=1$\,s later, at the read
$r(j)=\min(\{r: t_r\ge e_j+\epsilon\}\cup\{R\})$, and the target of read $r$ is the chunks due at $r$, serialised as in Eq.~\eqref{eq:serial}, then $\eta_r$; each read is
one training row whose prefix is the targets of all earlier reads. A deadline model is trained on these rows (new LoRA
adapters, the projector and the new rows) with Eq.~\eqref{eq:loss}, starting from \sys{} (offline). One such model,
trained on both directions (\S\ref{sec:setup}), is the Thinker of \sys{}; with the base Talker it is the
\emph{deadline baseline} (B), and it initialises write alignment (\S\ref{sec:write}). \sys{}-Simul, used in the
appendix analyses, has one deadline model per direction.

\subsection{Speaking in each speaker's voice}
\label{sec:method-voice}

\paragraph{Voices.} Runs of at least 0.4\,s in which slot $k$ alone is active (Eq.~\eqref{eq:alone},
Appendix~\ref{app:model}) are concatenated into a voice prompt $\bm{v}_k$ of at most 8\,s, in streaming mode only from
audio heard before the slot's first commitment. The UniSS Talker~\citep{uniss2025}, which predicts BiCodec
tokens~\citep{wang2025spark}, speaks each slot's segments with prompt $\bm{v}_k$. Each speaker has a stream of its own,
in which a segment starts at its release time or when the previous one ends (Eq.~\eqref{eq:layout}). The base Talker
plays every committed block in full at a fixed rate, so a stream that receives more speech than its source builds a
queue; a \emph{rule-based} Talker instead speeds up with the backlog\ifFactorial{} (Appendix~\ref{app:factorial})\fi.

\subsection{Playback alignment: when the translation is heard}
\label{sec:method-plem}

\begin{figure}[t]
\centering
\begin{tikzpicture}[x=0.82cm,y=0.6cm,font=\scriptsize,
  st/.style={circle,draw,inner sep=0pt,minimum size=4.2mm},
  on/.style={st,fill=trained!25},
  blk/.style={rectangle,draw,rounded corners=1pt,inner sep=1.5pt,minimum height=3.2mm}]
% time axis with release times
\draw[->,gray] (-0.3,-3.75) -- (7.4,-3.75) node[right,black]{$t$};
\foreach \x/\l in {0/\nu_i,2.4/\nu_{i+1},4.8/\nu_{i+2}}{\draw[gray] (\x,-3.65)--(\x,-3.85) node[below,black]{$\l$};}
% backlog states (rows = backlog values)
\foreach \c/\x in {i/0,i+1/2.4,i+2/4.8}{\node[gray] at (\x,1.05) {$h_{\c}$};}
\foreach \i/\x in {0/0,1/2.4,2/4.8}{\foreach \j in {0,1,2}{\node[st] (s\i\j) at (\x,-\j) {};}}
\node[left,gray] at (-0.25,0) {0\,s};\node[left,gray] at (-0.25,-1) {0.4\,s};\node[left,gray] at (-0.25,-2) {0.8\,s};
% candidate edges from h_i = 0
\draw[->,thick,trained] (s00) -- node[above,sloped]{$c$: short} (s10);
\draw[->,thick,frozen,dashed] (s00) -- node[below,sloped,pos=0.6]{long} (s12);
\draw[->,thin,gray] (s10) -- (s20); \draw[->,thin,gray] (s10) -- (s21);
\draw[->,thin,gray] (s12) -- (s21); \draw[->,thin,gray] (s12) -- (s22);
% backward message
\draw[->,densely dotted,thick,spkc] (5.6,-2.4) to[bend left=12] node[below,pos=0.45]{backward message $\beta_{i+1}(h)$} (0.6,-2.55);
\node[right,align=left] at (5.15,-0.35) {$\phi_i(h,c)=$\\$0.5\bar\ell(c)+5f(c)$\\$-\lambda L_i(h,c)$};
\end{tikzpicture}
\caption{Playback alignment for one speaker stream. States are the backlog $h_i$ (time block $i$ waits for earlier
speech); an edge is a candidate rendering $c$ of block $i$, whose duration sets the next backlog
(Eq.~\eqref{eq:backlog}). A long candidate delays every later block; backward messages bring that cost into the choice
for block $i$.}
\label{fig:pa}
\end{figure}


\paragraph{Lattice.} A stream's blocks $i=1,\dots,I$ have $W_i$ target words and are released at $\nu_i$
(\S\ref{sec:method-online}). For each block the Talker proposes candidates $c\in C_i$ (8 samples at two speaking rates)
of duration $D(c)$. The backlog $h_i$, how long block $i$ waits for earlier speech of its speaker, is rounded up
($\lceil\cdot\rceil$) to an 80\,ms grid up to 8\,s: $h_{i+1}=g_i(h_i,c_i)$ with
\begin{equation}
g_i(h,c)=\big\lceil\max(0,\,\nu_i+h+D(c)-\nu_{i+1})\big\rceil,
\label{eq:backlog}
\end{equation}
and with a block's words spread evenly over its playback they are heard $L_i(h,c)=W_ih+\tfrac{W_i+1}{2}D(c)$
word-seconds after $\nu_i$ in total. A linear-chain CRF over $(h_i,c_i)$ (Figure~\ref{fig:pa}) has potentials
\begin{equation}
\phi_i(h,c)=0.5\,\bar{\ell}(c)+5\,f(c)-\lambda L_i(h,c),
\label{eq:plpot}
\end{equation}
with $\bar{\ell}$ the Talker's mean token log-probability and $f$ the negative CTC loss per character under an
independent recogniser~\citep{baevski2020wav2vec}. A long candidate delays every later block of the stream, so the choice for block $i$
depends on the future; forward--backward over the backlog sums exactly over all candidate sequences \emph{within this
state space} (fixed candidates, the 80\,ms grid and the cap). The weights $0.5$ and $5$ were fixed before training and
are not ablated.

\paragraph{Budgets and training.} With $q_0$ the CRF at $\lambda=0$, $\lambda$ is the dual
variable of $\min_q\mathrm{KL}(q\Vert q_0)$ subject to $\mathbb{E}_q[\sum_iL_i]/\sum_iW_i\le b$, a posterior-regularised problem~\citep{ganchev2010posterior} whose solution is the CRF of Eq.~\eqref{eq:plpot}; the
expected delay falls monotonically in $\lambda$, so bisection finds $\lambda^\star$. The budget $b$ is the rule-based
Talker's own delay minus 0.3\,s. The Talker is then fine-tuned \emph{once} on each block's candidates, weighted by their
posterior marginal at the backlog the rule-based path reached, so the backward messages bring the queueing cost of
later blocks into the target for the current one (the \emph{playback-aligned}, PA, Talker; Appendix~\ref{app:plem}).
This is one round of posterior-weighted fine-tuning, not iterated EM, and no CRF runs at inference: the PA Talker,
conditioned on a backlog token, samples two candidates per block and keeps the one with the better CTC fidelity. The
budget therefore constrains the training posterior, not the online delay, which we measure. The \emph{playback layer}
(PA) is read-level merging (\S\ref{sec:method-online}) plus the PA Talker; \sys{} is the deadline Thinker with this
layer.

\subsection{Running on a whole meeting}
\label{sec:method-online}

\paragraph{Online windows.} A window opens at the first frame where any slot of the low-latency diarizer is active and
closes after 0.6\,s with no active slot (chosen on development meetings) or at 20\,s. All audio is processed. Commits
wait until the diarizer has released the frames that open the window (for the last read, those that close it), timed
as it emits them in blocks of nine (\emph{chunk protocol}, Appendix~\ref{app:chunk}). Slots persist across windows; the decoder restarts at each window. Blocks of one speaker
committed in the same read are merged before they reach the Talker (\emph{read-level merging}, rl).
